\section{Movement Phase}
During this phase, orders are revealed and detachments perform their movements.

\subsection{Activation}

Activating a detachment means revealing its order counter. The player with the initiative activates the first detachment.

Players then alternate, each activating one detachment at a time until no detachments remain to be activated. A player may choose to stop activating detachments. If they do so, they cannot activate any further detachments during the Movement Phase, except through reactions.

When a detachment is activated, any movement it performs is resolved immediately.

Off-table detachments that do not enter the battlefield cannot be activated. A detachment also cannot be activated if all its bases are pinned in an assault and are unable to perform an action.

\subsection{Movement}

\subsubsection{No Movement}

A detachment that cannot move, such as an immobilised detachment or one with a First Fire order, may still be activated during the Movement Phase. In this case, it cannot move or change its facing.

\subsubsection{Performing a Movement}

The maximum distance a base may move depends on its Movement characteristic and its order. When moving a base, measure the distance travelled by the centre of the base. A base may make any number of rotations during its movement, and these rotations do not reduce the distance it may move.

A base must be physically able to pass through every point along its movement path.

A base may only move through an enemy base if it is of a higher class, unless the enemy base is pinned in an assault. Immobilised bases of the same or a higher class cannot be moved through. Detachment coherency must be maintained throughout the entire movement.

The moving player must declare the route taken by the detachment, as this may be important for enemy reactions. If no route is declared, the detachment is considered to have taken the most direct route.

\paragraph{Reverse Movement}

A base may move in reverse. Each centimetre travelled in reverse counts as 2~cm of movement. Reverse movement allows a base to avoid exposing its vulnerable points and is only relevant to Class 3 and higher bases.

\paragraph{Movement and Zones of Control}

A base may only move through an enemy base's zone of control if it is of a higher class than the enemy base. Engaging in an assault is an exception to this restriction.

If a base begins its movement within a zone of control, that zone of control is ignored for the duration of that movement.

A base is considered to be within another base's zone of control if any part of it is inside the zone, not only its centre.

\paragraph{Engaging in an Assault}

For one of your detachments to engage enemy bases in an assault, it must have a Charge order and sufficient movement to make contact with its target. A base is engaged in an assault if it is in contact with an enemy base. A base that engages in an assault loses its ability to fire during that turn.

A base may move through the zones of control of enemy bases that it intends to engage without restriction. It may also move through multiple overlapping zones of control to make contact with its target.

A base must always face the enemy when engaging in an assault. When it makes contact with an enemy base, that enemy must therefore be within its front arc. Consequently, a base cannot engage in an assault while moving in reverse.

\begin{figure}
    \centering
    \includegraphics[width=0.5\linewidth]{movement-a.png}
    \caption{The assault engagement is not valid in the second example because the base is not engaging the enemy within its front arc.}
    \label{fig:MovementA}
\end{figure}

All bases in an enemy detachment must be engaged in an assault before the attacking detachment may engage another enemy detachment. However, the attacker may ignore bases that it cannot engage—for example, if it lacks sufficient movement to contact an enemy base or if there is not enough room to place its own base in contact.

These engagement rules apply both to the detachment initiating the assault and to any reinforcements that subsequently join it.

\begin{figure}
    \centering
    \includegraphics[width=0.5\linewidth]{movement-b.png}
    \caption{The red detachment must fully engage the yellow detachment before it can engage the blue detachment.}
    \label{fig:MovementB}
\end{figure}

Each base may normally engage only one enemy base in an assault. However, if a base has a higher class than all of its opponents, it may engage multiple enemy bases of a lower class simultaneously.

In this case, it may engage one additional enemy base for each point of difference between its class and the highest class among its opponents.

\begin{figure}
    \centering
    \includegraphics[width=0.5\linewidth]{movement-c.png}
    \caption{A Class 3 Haruspex may engage two Class 2 Space Marine Bike bases (Example 1) or three Class 1 Tactical Space Marine bases (Example 2). It may also engage one Space Marine Bike base and one Tactical Space Marine base, but no more. Because the two opposing bases have different classes, the Cavalry limitation applies (Example 3).}
    \label{fig:MovementC}
\end{figure}

\paragraph{Pushing Bases}

A base may be positioned where enemy bases of a lower class are already located, unless those enemy bases are already engaged in an assault. The owner of the pushed bases must reposition them as close as possible to their original positions. This may cause their detachment to lose coherency.

Bases may be pushed while simultaneously being engaged in an assault, provided that the standard requirements for engaging in an assault are met.

If the number of pushed bases is equal to or lower than the number of bases that may be engaged in the assault, all pushed bases are engaged. If more bases are pushed than can be engaged, any bases that are not engaged are placed 1~cm from the enemy base without being engaged in the assault. These bases may still consolidate as a reaction.

\begin{figure}
    \centering
    \includegraphics[width=0.5\linewidth]{movement-d.png}
    \caption{The Class 4 Barbed Hierodule engages a detachment of six Class 1 Tactical Space Marine bases. The Hierodule engages and pins four bases. The other two bases are not engaged and are moved a few millimetres aside.}
    \label{fig:MovementD}
\end{figure}

\paragraph{Disengaging from an Assault}

If a base is engaged in an assault but is not pinned, it may disengage from that assault. Disengaging can be dangerous.

When a base disengages, each enemy base in contact with it may attempt an opportunity attack. The result required to inflict a hit depends on the enemy base's Assault Factor (AF), as shown in the following table:

\begin{table}[ht]
    \centering
    \begin{tabular}{|l|c|c|c|}
        \hline
        \textbf{AF of the Base Making the Opportunity Attack}
            & \textbf{Less than 7}
            & \textbf{7--12}
            & \textbf{13+} \\
        \hline
        To-Hit Roll
            & 6+
            & 5+
            & 4+ \\
        \hline
    \end{tabular}
\end{table}

Hits inflicted in this manner are treated as assault hits and therefore allow no armour save. They retain all their normal characteristics, including any assault damage bonuses.

\begin{figure}
    \centering
    \includegraphics[width=0.5\linewidth]{movement-e.png}
    \caption{The Barbed Hierodule has been engaged by four Tactical Space Marine bases (AF 2) and one Assault Terminator base (AF 7) and decides to disengage.    
    It therefore suffers four opportunity attacks hitting on 6+ from the Tactical Space Marines and one hitting on 5+ from the Assault Terminators.}
    \label{fig:MovementE}
\end{figure}


\subsection{Reactions to Movement}

A player may react to enemy movement in three main ways, as described below. These reactions do not count as activations during the Movement Phase.

\subsubsection{Overwatch Fire}

If an enemy detachment moves within the line of sight and range of one of your detachments with an unrevealed First Fire order, you may perform Overwatch Fire.

Several of your detachments may perform Overwatch Fire at different points during the enemy detachment's movement. To do so, pause the enemy detachment's movement at the point where you wish to fire and resolve the Overwatch Fire. The enemy detachment may then continue its movement, provided it is still able to do so.

A detachment that performs Overwatch Fire cannot fire or be activated again for the remainder of the turn, and its order counter is removed. The entire detachment is activated for the Overwatch Fire, even if some of its bases are unable to fire. Detachments performing Overwatch Fire may reroll results of 1 on their To-Hit rolls.

Overwatch Fire is resolved in the same way as normal shooting, including rerolling results of 1, but suffers a $-1$ penalty to its To-Hit rolls unless:

\begin{itemize}
    \item The targeted bases are within the zone of control of at least one base belonging to the firing detachment.
    \item The firing detachment is being charged by the targeted detachment.
\end{itemize}

The player receiving the fire chooses which bases in the targeted detachment are hit. Casualties are removed immediately.

Ranged template weapons cannot perform Overwatch Fire. Rangeless template weapons may perform Overwatch Fire. When doing so, each rangeless template weapon rolls three attack dice and has a range of 20~cm.

If damage caused by Overwatch Fire leaves the targeted detachment without sufficient movement to reach its intended target, it stops after moving as far as possible. It cannot engage bases other than its original target in an assault.

\subsubsection{Consolidation}

If an enemy base charges or ends its movement within the zone of control of one or more bases from one of your detachments, every unengaged base in that detachment may immediately make a 6~cm \textit{Consolidation} move after the enemy finishes moving.

This move may allow the bases to engage enemy bases within 6~cm in an assault. Bases that are too far away to engage or join an assault may also consolidate, but only towards the enemy bases that triggered the consolidation and only to maintain their detachment's coherency.

A detachment cannot consolidate if it has a Fall Back order, is immobilised, or has performed Overwatch Fire.

A Consolidation move may be used to leave a structure or transport. It cannot be used to disengage, enter a garrison, or board a transport.

If one of your bases consolidates, it is considered to have engaged in an assault, and all corresponding rules apply, including losing its zone of control and its ability to fire.

\subsubsection{Countercharge}

To perform a Countercharge, a detachment must have been activated during the Movement Phase and placed in a Countercharge position. A base's Countercharge radius is equal to its Movement characteristic.

If an enemy detachment ends its movement within this radius, the detachment may move to engage it in an assault. This is a standard movement that follows the normal movement rules, except that it is performed as a reaction to an enemy activation.

A Countercharge cannot be used to bring reserves onto the battlefield or to disengage from an assault.